\documentclass{article}
\usepackage[T1]{fontenc}
\usepackage{lmodern}
\usepackage{graphicx}
\usepackage{amsmath,amssymb}
\usepackage[a4paper,margin=2cm]{geometry}
\usepackage[absolute,overlay]{textpos}
\setlength{\TPHorizModule}{1mm}
\setlength{\TPVertModule}{1mm}
\usepackage[indonesian]{babel}
\usepackage{hyperref}
\setlength{\emergencystretch}{2em}
% unit: Halaman 13

\begin{document}

% === Halaman 13 (python/native) ===

13

Contoh 1: Diberikan cipherteks berikut ini (Stalling, 2011), spasi tidak 

dibuang:

 UZ QSO VUOHXMOPV GPOZPEVSG ZWSZ OPFPESX UDBMETSX AIZ 

VUEPHZ 

HMDZSHZO 

WSFP 

APPD 

TSVP 

QUZW 

YMXUZUHSX 

EPYEPOPDZSZUFPO MB ZWP FUPZ HMDJ UD TMOHMQ

 Kita akan malakukan kriptanalisis dengan metode analisis frekuensi untuk 

memperoleh plainteks. 

   Asumsi: bahasa yang digunakan adalah Bahasa Inggris dan cipher yang digunakan 

adalah cipher abjad-tunggal.

\end{document}
